\section{Related Work}
\label{sec:related}

\subsection{Agent Skills as a Supply-Chain Boundary}

Agent Skills occupy an unusual point between prompts and packages: they combine
natural-language control with executable code, configuration, and resources.
Liu et al. establish the ecosystem threat through a large-scale measurement of
public registries, using static candidate generation followed by sandbox
execution and expert confirmation~\cite{liu2026maliciousskills}. MalSkills
instead targets zero-execution detection by extracting security-sensitive
operations, linking values across artifacts, and applying neuro-symbolic rules
to a dependency graph~\cite{wang2026malskills}. MaliciousSkillBench consolidates
multiple Skill sources and demonstrates that performance degrades when the
evaluation separates provenance~\cite{wang2026maliciousskillbench}. Together,
these works define the central tension for admission control: behavioral
confirmation requires execution, whereas artifact-only detection must reason
under ambiguity and distribution shift.

The broader software-supply-chain literature exposes a complementary entry
point. Code-generating models frequently hallucinate plausible package names,
creating opportunities for package-confusion attacks~\cite{spracklen2025package}.
That work asks which dependency an LLM causes a developer to select; malicious
Skill detection asks what an already selected package can instruct and enable
an agent to do. Our work is closest to MalSkills, but evaluates whether adding
separated high-level context to recovered behavior improves detection under a
source-disjoint protocol. The result is negative, delimiting this design point
rather than treating richer structure as an assumed benefit.

\subsection{Adversarial Instructions, Tools, and Agent State}

Prompt injection studies an attacker who places competing instructions in
content consumed by an application. Liu et al. formalize this problem and
benchmark attacks and defenses across models and tasks
~\cite{liu2024promptinjection}. InjecAgent specializes the threat to
tool-integrated agents, covering harmful actions and private-data exfiltration
triggered by tool-returned content~\cite{zhan2024injecagent}. These attacks occur
during a task trajectory, whereas a malicious Skill can make adversarial
instructions persistent and package them with helper code or installation
behavior.

Other work compromises different components of the agent stack. AgentPoison
plants retrieval triggers in long-term memory or knowledge bases
~\cite{chen2024agentpoison}; BadAgent inserts backdoors through an untrusted
model or fine-tuning data~\cite{wang2024badagent}. ToolSword maps failures across
tool-use input, execution, and output stages~\cite{ye2024toolsword}, while
ToolAlign trains for helpful, harmless, and autonomous tool use
~\cite{chen2024toolalign}. These studies show that agent behavior can be
subverted through model weights, memory, retrieved content, or tool feedback.
Our adversary instead controls a distributable Skill artifact. The detector
therefore cannot retrain the underlying model or observe a trigger at runtime;
it must identify suspicious instructions and implementation evidence before
the artifact becomes agent state.

\subsection{Dynamic Analysis and Agent Security Evaluation}

Dynamic approaches judge security from execution trajectories. AgentFuzz uses
directed greybox fuzzing, generated seeds, and semantic and distance feedback
to trigger taint-style vulnerabilities in running agent applications
~\cite{liu2025agentfuzz}. ToolEmu searches for long-tail failures in
language-model-emulated tool environments without invoking real services
~\cite{ruan2024toolemu}. AgentDojo combines utility tasks, security goals, and
indirect prompt-injection attacks in a reproducible agent environment
~\cite{debenedetti2024agentdojo}; Agent Security Bench spans prompt, memory,
tool, and multi-agent attack surfaces~\cite{zhang2025asb}.

These evaluations observe whether an agent performs a prohibited action, which
provides a semantically direct outcome but depends on an executable system and
a triggering interaction. Static package analysis moves the decision earlier:
it can cover dormant artifacts without running them, but must expose uncertain
reachability, unsupported languages, and incomplete parsing. Our evaluation
therefore reports class-conditional detection together with processing
coverage, abstention, and unresolved outputs. It does not treat the absence of
a recovered static path as evidence that no runtime behavior exists.

\subsection{Instruction Integrity and Runtime Enforcement}

Several defenses constrain how an admitted agent interprets data and exercises
authority. StruQ separates instructions from data through a structured query
interface and model training~\cite{chen2025struq}. Task Shield checks whether
candidate instructions and tool calls advance the user's declared objective
~\cite{jia2025taskshield}. CaMeL derives control flow from trusted input and
tracks capabilities on values passed to tools~\cite{debenedetti2026camel}.
These mechanisms target indirect injection during task execution; they can
mediate actions even when the underlying package is already present.

System mechanisms enforce the same principle below the prompt layer.
IsolateGPT places third-party agent applications in separate contexts with
explicit interfaces~\cite{wu2025isolategpt}, while AgentSpec provides a policy
language for runtime constraints over agent actions~\cite{wang2026agentspec}.
Isolation and policy enforcement limit consequences but do not decide whether
an artifact is trustworthy before installation. Conversely, a static
admission decision cannot observe environment-dependent behavior and cannot
replace least privilege after admission. Read Before You Run occupies this
upstream boundary: it produces source-grounded pass/review/block evidence for
the package, leaving runtime mediation to complementary defenses.
